\section{Why the Responses Vary: Conditions and Mechanism}
\label{sec:why}

Section~\ref{sec:observed} established what a law must reproduce. We now examine which response structures can be shared across capabilities, how model state and training protocol affect prediction, and how logit displacement accounts for compression damage.


\subsection{Function shape}
\label{sec:why_shape}

Under pruning, once the measured first-order alignment term is subtracted, the residual before the
cliff is a power law in the deleted capability mass, with an exponent that transfers across
families; this description uses gradient information from the dense
model, a richer input than \S\ref{sec:laws} admits. Under quantization the response is instead a
threshold in bit-width whose position depends on the group size, so an interpolation over the
measured grid predicts it where a smooth surface does not.

Under distillation, a joint form with a logarithmic reuse term,
$\delta=(a+a'z)\log(1+T/T_{\mathrm{ref}})+(b+b'z)\log(1+E)$ with $z$ the standardized student size
and $T_{\mathrm{ref}}$ a reference budget, implies that the pool-change effect, divided by its logarithmic bracket, is flat in budget. It is
not: for code on the 1B student it drifts by a factor of twenty across checkpoints (Appendix
Fig.~\ref{fig:drift}). A free reuse exponent, profiled inside training folds, lies near or above one for every capability,
though with wide fold ranges (Appendix Fig.~\ref{fig:explanation}a). Since neither the logarithmic
form, the curved form nor a single interaction term improves on the strongest development-stage
baseline at either pre-declared target, we deliver the reuse term in its logarithmic form.

Those data cannot say whether additivity itself failed, because a fixed pool and a doubling budget
grid never place two trajectories at the same budget. We therefore designed four trajectories by
replaying the update boundaries of the trainer, and they landed on a rectangle to the token. The second difference across that rectangle, which is zero for any additive structure whatever its
curvature, was $-0.24$ nats on the 1B student and $-0.21$ on the 4B for the question-answering
probe. Both lie inside the pre-registered band of $\pm0.32$, and both fall between the additive
prediction of zero and the $-0.49$ of the frozen interaction candidate
(Appendix Fig.~\ref{fig:corner_test}). The pre-fixed rule neither rejects nor establishes additivity; the 4B statistic on mathematics and code
fell outside its band at readouts that were registered as underpowered in advance.



\subsection{Source state, protocol, and evaluation distribution}
\label{sec:why_conditions}

Size and stage carry mathematics and code to unseen densities of states that entered the fit, and
they over-predict question answering, whose density response is not proportional to any shared
curve. Dense pre-compression statistics are less accurate than the source-free median curve in nine
of nine pairs. Under distillation the learning rate sets the magnitude of the damage, and at a fixed rate the
reuse ratio $E=T/D_U$ dominates the factors we varied. The question-answering cost it sets is
monotone, grows with student size and transfers across the three distributions, while the gain is
confined to 2Wiki; a residual dependence on budget, pool composition and student remains. A data requirement phrased as a bound on the loss increase is therefore supportable across
distributions, whereas a claim about gain holds only for 2Wiki among the three we scored.

\subsection{Mechanism}
\label{sec:why_mech}

For pruning and quantization, the loss change under a logit displacement $r$ obeys the identity
$\Delta L=(\mathbb{E}_{\pi}[r]-r_y)+\tfrac12\mathrm{Var}_{\pi}(r)+O(\lVert r\rVert^3)$, where $\pi$ is
the next-token distribution of the reference model and $r_y$ the displacement at the target token,
and its coefficient-free truncation, the mean displacement against the target token plus half the
displacement variance, reproduces the measured damage under pruning, grouped round-to-nearest and
per-channel round-to-nearest quantization with an accuracy that follows the damage magnitude on the
measured panel: under 5\% median relative error below about 0.1 nats and past 20\% by 0.25
(Appendix Fig.~\ref{fig:explanation}b). Both terms carry the account; the first-order term is not
negligible, and the two can cancel. The range is empirical, since the remainder is controlled by the size of the displacement. Pure logit shrinkage carries about three percent of the response, and the residual variance that
carries the rest needs the compressed model, so the account explains damage after the fact and
enters no relation as an input.
